% Preamble
\documentclass[a4paper,12pt]{article}

% Prevents line breaks mid-word
\tolerance=1
\emergencystretch=\maxdimen
\hyphenpenalty=10000
\hbadness=10000

% Packages
\usepackage{amsfonts}
\usepackage{blindtext}
\usepackage{titlesec}
\usepackage{todonotes}
\usepackage{graphicx}
\graphicspath{{images/}}
\usepackage[hidelinks]{hyperref}

 % Set inline as default for todos
\usepackage{regexpatch}
\makeatletter
\xpatchcmd{\@todo}{\setkeys{todonotes}{#1}}{\setkeys{todonotes}{inline,#1}}{}{}
\makeatother

% Document
\begin{document}
    \begin{titlepage}
        \centering
        \includegraphics[width = 1\textwidth]{UoB_Logo}\\[10ex]
        \Huge{A Comparison of Approaches to Combinatorial Optimisation for Multi-Day Route Planning}\\[14ex]
        \LARGE{Jacob Luck}\\
        \large{2338114}\\[2ex]
        \Large{B.Sc. Computer Science with a Year in Industry}\\
        \Large{Project Supervisor - Leandro Minku}\\[2ex]
        \large{April 2025}
    \end{titlepage}

    \section*{Abstract}\label{sec:abstract}
    \todo{Write abstract}

    \pagebreak

    \tableofcontents

    \pagebreak

    \section{Introduction}\label{sec:introduction}
    \subsection{Motivation}\label{subsec:motivation}
    \todo{Write about reasoning for this project, can copy a little bit over from presentation slides. Explain problem informally. Link into aims and objectives.}
    \subsection{Aims and Objectives}\label{subsec:aims-and-objectives}
    \todo{Explin goals of the project, link in to methodology.}
    \subsection{Methodology}\label{subsec:methodology}
    \todo{Briefly explain how the project will be carried out. A more thorough description can be provided later on, i.e., when describing algorithms.}
    \subsection{Summary}\label{subsec:summary}
    \todo{Explain what is in the rest of the report. "In this report I shall...", cover each section, etc.}

    \pagebreak
    
    \section{Problem Formulation}\label{sec:problem-formulation}
    \subsection{Problem Description}\label{subsec:problem-description}
Given a positive integer $d$ and a graph $G = (V, E)$, where $V$ is set of locations including a designated
starting point $s$ and $E$ is a set of weighted edges linking every location to every other location, find a route
that:
\begin{enumerate}
    \item Visits all nodes $V \setminus c$ once.
    \item Starts and finishes at $s$, having visited it $d$ times, without ever visiting consecutively.
    \item Minimises both the cumulative edge weights in the route and the variance in cumulative weight between
    each visit to $s$.
\end{enumerate}

\subsection{Inputs, Outputs and Design Variables}\label{subsec:inputs-and-outputs}
Inputs:
\begin{itemize}
    \item $d$: The number of times $s$ should be visited in a route.
    Contextually, $d$ represents the number of days a tourist will spend on their trip. $d \in \mathbb{Z}, d > 0$
    \item $G = (V, E)$: A pair comprising:
    \begin{itemize}
        \item[\textbullet] $V$: A set of nodes representing locations the tourist would like to visit.
        $v \in V, v = (x, y, t)$, a triple comprising:
        \begin{itemize}
            \item[\textbullet]$x$: Longitude, indicating the location's geographic east-west position on the
            earth\footnote{\label{latitude-longitude-note}While the coordinates of our locations are
            included in $V$, they are not directly tied to the weight of our edges $E$, which are based on
            time and not distance.}.\\
            $x \in \mathbb{Q}, -180 \leq x \leq 180$.
            \item[\textbullet]$y$: Latitude, indicating the location's geographic north-south position on
            the earth\footnote{See footnote \ref{latitude-longitude-note}}.\\
            $y \in \mathbb{Q}, -90 \leq y \leq 90$.
            \item[\textbullet]$t$: Duration, in minutes, indicating how much time to spend at this location.\\
            $t \in \mathbb{Z}, t > 0$.
        \end{itemize}
        \item[\textbullet] $E$: A set of edges $e \in E$ that connects every node to every other node,
        bidirectionally. $e = (v_1, v_2, w)$, a triple comprising of:
        \begin{itemize}
            \item[\textbullet]$v_1$: A location representing the origin of the edge.\\
            $v_1 \in V$.
            \item[\textbullet]$v_2$: A location representing the destination of the edge.\\
            $v_2 \in V$.
            \item[\textbullet]$w$: A weight indicating the sum of the time it takes to travel from $v_1$ to $v_2$
            and the time the tourist wishes to spent at $v_2$.\\
            $w \in Z, w > 0$.
        \end{itemize}
    \end{itemize}
    \item $s$: Starting point that should be visited $d$ times.
    Contextually, $s$ represents where the tourist is staying and will return to at the end of each day.\\
    $s \in V$.
\end{itemize}
Outputs:
    \begin{itemize}
    \item $R$: A valid route satisfying all constraints, represented as an ordered sequence of locations.\\
    $R = [r_1, r_2, \dots, r_n], r_i \in V$.
\end{itemize}

\subsection{Objective Function}\label{subsec:objective-function}
As previously mentioned in the \hyperref[subsec:problem-description]{Problem Description}, our goal is to find a
route that minimises the cumulative weight and the variance in route weight between each visit to $s$.
To accomplish this the following cost function is applied to each route:
\todo{Update equation for cost, talk to leandro about how to handle changing this.}
\begin{equation}
    Cost(R) = W \times (1 + \sigma)\label{eq:cost}
\end{equation}
Where $W$ is the sum of the weights of all edges traversed in the route and $\sigma$ is the standard deviation of the
sum of weights between each visit to $s$:
\begin{equation}
    W = \sum_{i=0}^{n-1} w(r_i, r_{i+1}), r_i \in R\label{eq:weight}
\end{equation}
Where $w(r_i, r_{i+1})$ is the weight of the edge between $r_i$ and $r_{i+1}$.
\begin{equation}
    \sigma = \sqrt{\frac{\sum_{i=0}^{d}(x_i-\mu)}{d}}, x_i \in X\label{eq:standard-deviation}
\end{equation}
Where $R$ is divided into sections between each visit to $s$ and $X$ is a list of the sum of weights within these
sections.\\
$\mu$ is the mean cumulative weight of each $x_i$.

\subsection{Constraints}\label{subsec:constraints}
\todo{Describe constraints}

    \pagebreak

    \section{Literature Review}\label{sec:literature-review}
    \todo{Plan (and write) literature review}
    \todo{Remember to write about the strengths and weaknesses of existing work. At the end of this chapter you can then give a summary of the gaps that you'll be trying to improve with your work, and on the strengths that you will be maintaining in your work.}
    This literature review aims to explore existing research and approaches to other combinatorial optimization problems.
There is extensive previous research on various combinatorial problems, for example, the Travelling Salesman
Problem (TSP), Vehicle Routing Problem (VRP) and Tourist Trip Design Problem (TTDP).
It is important to understand how these problems are similar to the one presented in this report, as well as where
those similarities end.
By gaining an understanding of the strengths and limitations of existing approaches to similar problem, we can make
better informed decisions regarding which approaches to investigate, how they may be adapted to suit our specific
constraints and how they might be implemented in practice.
While the approaches taken to these problems may not be directly applicable to our own, it is likely we can adapt their
techniques to suit the specific constraints of this problem.

[paper1] defines [problem1] as [definition].
blah blah blah.
This is similar to our problem because [reason].
Another problem is [problem1], which is defined by [paper2] as [definition].
blah blah blah.
This is similar to our problem because [reason].

[paper3] proposes a solution to [problem] using [approach].
blah blah blah, pros, cons, etc.

READ, YOU NEED TO READ THE PAPERS stop faffing about and read the papers


\subsection{Related Optimization Problems}
\begin{itemize}
    \item Vehicle Routing Problems (VRP), particularly Multiple-Trip VRP (Cattaruzza et al., 2016)
    \item Traveling Salesman Problem variants (Applegate et al., 2006)
    \item Tourist Trip Design Problem (TTDP) (Vansteenwegen et al., 2011; Gavalas et al., 2014)
    \item Orienteering Problems with time constraints (Gunawan et al., 2016)
\end{itemize}

\subsection{Algorithms for Multi-Visit Routing}
\begin{itemize}
    \item Exact methods: Branch and Bound, Dynamic Programming, Integer Linear Programming (Laporte, 1992)
    \item Heuristic approaches for multi-visit scenarios (Necula et al., 2015)
    \item Metaheuristics: Genetic Algorithms (Potvin, 1996), Simulated Annealing (Černý, 1985), Ant Colony Optimization (Dorigo \& Gambardella, 1997)
    \item Approaches handling returning to a depot multiple times (Azi et al., 2014)
\end{itemize}

\subsection{Multi-Objective Optimization Approaches}
\begin{itemize}
    \item Survey of approaches for handling multiple objectives
    \begin{itemize}
        \item Weighted sum methods (Marler \& Arora, 2010)
        \item Pareto optimization approaches (García-Nájera et al., 2013)
        \item Constraint-based methods (Konak et al., 2006)
    \end{itemize}
    \item Techniques for minimizing variance/balancing workloads (Dutot et al., 2006; Ghannadpour et al., 2013)
\end{itemize}

\subsection{Time-Constrained Route Optimization}
\begin{itemize}
    \item Handling time windows and duration constraints (Toth \& Vigo, 2002)
    \item Techniques for time-dependent weights (Gendreau et al., 2015)
    \item Applications to tourist trip planning with time budgets (Souffriau et al., 2013)
\end{itemize}

\subsection{Real-World Applications}
\begin{itemize}
    \item Implementation considerations for tourist route planning systems (Souffriau et al., 2008)
    \item Personalized itinerary systems and their algorithms (Schaller et al., 2014)
    \item Case studies of similar optimization problems in practice (Gavalas et al., 2015)
\end{itemize}

\subsection{Gap Analysis and Research Directions}
\begin{itemize}
    \item Identify how existing approaches fall short for this specific combination of constraints
    \item How methods might be adapted or combined to address the problem (Liao et al., 2017)
    \item Novel aspects requiring new algorithmic development (Nazari et al., 2018; Chou et al., 2020)
\end{itemize}

\subsection{Conclusion}
\begin{itemize}
    \item Summarize the key algorithmic approaches most relevant to this problem
    \item Identify the most promising directions based on the literature
    \item Transition to the approach taken in the current research
\end{itemize}

Full references (you'll need to format these according to your preferred citation style):

Applegate, D. L., Bixby, R. E., Chvátal, V., & Cook, W. J. (2006). The traveling salesman problem: a computational study. Princeton University Press.
Azi, N., Gendreau, M., & Potvin, J. Y. (2014). An exact algorithm for a vehicle routing problem with time windows and multiple use of vehicles. European Journal of Operational Research, 202(3), 756-763.
Cattaruzza, D., Absi, N., Feillet, D., & Vidal, T. (2016). The multi-trip vehicle routing problem with time windows and release dates. Transportation Research Part E: Logistics and Transportation Review, 92, 118-133.
Černý, V. (1985). Thermodynamical approach to the traveling salesman problem: An efficient simulation algorithm. Journal of Optimization Theory and Applications, 45(1), 41-51.
Chou, C. C., Fu, J. S., Kuo, S. P., & Li, P. H. (2020). A hybrid algorithm of simulated annealing and tabu search for the bi-objective team orienteering problem with time windows. IEEE Access, 8, 72564-72578.
Dorigo, M., & Gambardella, L. M. (1997). Ant colony system: a cooperative learning approach to the traveling salesman problem. IEEE Transactions on evolutionary computation, 1(1), 53-66.
Dutot, P. F., Laugier, A., & Bustos, A. M. (2006). Balancing a dynamic mobile workforce via constraint-based local search. In International Conference on Innovative Techniques and Applications of Artificial Intelligence (pp. 35-48).
García-Nájera, A., Bullinaria, J. A., & Gutiérrez-Andrade, M. A. (2013). A Pareto ant colony optimization algorithm for the multi-objective tourist trip design problem. In Applications of Evolutionary Computation (pp. 143-154).
Gavalas, D., Konstantopoulos, C., Mastakas, K., & Pantziou, G. (2014). Models and algorithms for the tourist trip design problem. Annals of Operations Research, 224(1), 65-86.
Gavalas, D., Konstantopoulos, C., Mastakas, K., & Pantziou, G. (2015). A variable neighborhood search approach for the tourist trip design problem. Expert Systems with Applications, 42(21), 7911-7919.
Gendreau, M., Ghiani, G., & Guerriero, E. (2015). Time-dependent routing problems: A review. Computers & Operations Research, 64, 189-197.
Ghannadpour, S. F., Noori, S., Tavakkoli-Moghaddam, R., & Ghoseiri, K. (2013). A multi-objective vehicle routing problem with soft time windows: The minimization of drivers' dissatisfaction. Transportation Research Part E: Logistics and Transportation Review, 53, 64-77.
Gunawan, A., Lau, H. C., & Vansteenwegen, P. (2016). Orienteering problem: A survey of recent variants, solution approaches and applications. European Journal of Operational Research, 255(2), 315-332.
Konak, A., Coit, D. W., & Smith, A. E. (2006). Multi-objective optimization using genetic algorithms: A tutorial. Reliability Engineering & System Safety, 91(9), 992-1007.
Laporte, G. (1992). The traveling salesman problem: An overview of exact and approximate algorithms. European Journal of Operational Research, 59(2), 231-247.
Liao, T. W., Egbelu, P. J., & Chang, P. C. (2017). Integrated manufacturing systems design with robot movement simulation and multi-objective genetic algorithms. International Journal of Production Research, 55(15), 4399-4422.
Marler, R. T., & Arora, J. S. (2010). The weighted sum method for multi-objective optimization: new insights. Structural and Multidisciplinary Optimization, 41(6), 853-862.
Nazari, M., Oroojlooy, A., Snyder, L., & Takác, M. (2018). Reinforcement learning for solving the vehicle routing problem. Advances in Neural Information Processing Systems, 31.
Necula, R., Breaban, M., & Raschip, M. (2015). Exact and heuristic approaches for the multi-visit salesman problem. In International Conference on Computational Collective Intelligence (pp. 357-368).
Potvin, J. Y. (1996). Genetic algorithms for the traveling salesman problem. Annals of Operations Research, 63(3), 337-370.
Schaller, R., Elsweiler, D., & Harvey, M. (2014). City trip planner: Personalized itinerary planning for groups of people. In Proceedings of the 22nd International Conference on User Modeling, Adaptation, and Personalization (pp. 184-191).
Souffriau, W., Vansteenwegen, P., Vertommen, J., Berghe, G. V., & Oudheusden, D. V. (2008). A personalized tourist trip design algorithm for mobile tourist guides. Applied Artificial Intelligence, 22(10), 964-985.
Souffriau, W., Vansteenwegen, P., & Oudheusden, D. V. (2013). The multi-constraint team orienteering problem with multiple time windows. Transportation Science, 47(1), 53-63.
Toth, P., & Vigo, D. (2002). The vehicle routing problem. SIAM.
Vansteenwegen, P., Souffriau, W., & Van Oudheusden, D. (2011). The orienteering problem: A survey. European Journal of Operational Research, 209(1), 1-10.

    \pagebreak

    \section{Algorithms Investigated}\label{sec:algorithms-investigated}
    \todo{Paragraph describing different types of algorithm used (Routing then cluster, Cluster then Routing, Genetic, etc.)}
    \todo{Remember to justify the choice of algorithms. You may also need to explain how to adopt these algorithms in your work. A figure showing the ralationship between different components of your work may also help.}
    \subsection{Routing}\label{subsec:routing}
    \todo{Explain purpose of routing/goal of algorithms.}
    \subsubsection{Brute Force}\label{subsubsec:brute-force-routing}
    \todo{Write brute force explanation}
    \subsubsection{Greedy Routing/Insertion}\label{subsubsec:greedy-routing}
    \todo{Explain greedy routing algorithm}
    \todo{Explain greedy insertion algorithm}
    \subsubsection{Gift Wrapping}\label{subsubsec:gift-wrapping}
    \todo{Explain gift wrapping algorithm}
    \todo{Something like: "Once gift wrapping has found a convex hull, a greedy insertion algorithm is used to find the optimal route within the convex hull."}

    \subsection{Clustering}\label{subsec:clustering}
    \todo{Explain purpose of clustering, how it is used in route planning and the goal of our algorithms.}
    \subsubsection{K-Means}\label{subsubsec:k-means}
    \todo{Explain k-means algorithm}

    \subsection{Trip Generation}\label{subsec:trip-generation}
    \todo{Explain trip generation, how it is used in route planning and the goal of our algorithms.}
    \subsubsection{Brute Force}\label{subsubsec:brute-force-trip-generation}
    \todo{Explain how brute force algorithm can be modified for trip generation.}

    \subsection{Genetic Algorithms}\label{subsec:genetic-algorithm}
    \todo{Explain basics of genetic algorithms, include how by considering different genomes, the algorithm can be used for routing, clustering or trip generation.}
    \subsubsection{General Clustering}
    \todo{Explain genetic clustering}
    \subsubsection{Centroid-based Clustering}
    \todo{Explain genetic centroid-based clustering and how it differs from general clustering.}
    \subsubsection{Routing}
    \todo{Explain genetic routing}
    \subsubsection{Trip Generation}
    \todo{Explain genetic trip generation}

    \pagebreak

    \section{Implementation}\label{sec:implementation}

    \subsection{Constraints}\label{subsec:implementation-constraints}

    \pagebreak

    \section{Evaluation and Comparison}\label{sec:evaluation-and-comparison}
    \todo{Write paragraph about experiment process. Comparison based on computation time and route evaluation.
    Describe how route is evaluated. Describe data being tested on.}
    \todo{Present comparison of different combinations of algorithms on different inputs.}
    \todo{Reflect about the questionns you are trying to answer with your evaluation. You can have onne subsection for each question that you are trying to answer. It's also important to justify your choices when it comes to the methodology used for evaluation.}

    \pagebreak

    \section{Conclusion and Future Work}\label{sec:conclusion-and-future-work}
    \todo{Write conclusion, discuss results, comparison of algorithms, etc.}
    \subsection{Project Reflection}
    \todo{Reflect on the project, what went well, what didn't go well, what I would do differently. This should lead into future work.}
    \subsection{Future work}\label{subsec:future-work}
    \todo{Discuss further work, what I will be doing to improve the project}

    \pagebreak

    \section{Bibliography}\label{sec:bibliography}
    \todo{Fill in bibliography}

\end{document}